\setcounter{section}{26}
\section{Protección de datos}

Nombre completo: La política de protección de datos de carácter personal. Régimen Jurídico. El reglamentoUE 2016/679, de 27 de abril, relativo a la protección de las personas físicas en lo que respecta al tratamiento de datos personales y a la libre circulación de estos datos. Principios y derechos. Obligaciones. El delegado de protección de datos en las administraciones públicas. La Agencia Española de Protección de Datos.

\localtableofcontents

\subsection{Contexto normativo}
Tradicionalmente se aplicaba la LO 5/1992 del tratamiento automatizado de datos peronales, entre otras normas. Actualmentescode tgex se aplica el Reglamento 2016/679.


\subsection{Política de protección de datos}
Una política de protección de datos da a conocer el modo en el que una empresa u organismo obtiene, trata y protege los datos personales de sus empleados, clientes o usuarios. También informa de los medios puestos a disposición para permitir ejercer los derechos a los propietarios de los datos.

\subsection{RGPD}
    Gira entorno al concepto de Registro de Tratamiento. Quedando reflejado así en su artículo 30. La protección de las personas físicas en relación al tratamiento de los datos personales es un derecho fundamental en la CDFUE y en el Tratado de Funcionamiento de la Unión Europea. 

    \subsubsection{Régimen Jurídico}
        Entró en vigor en mayo de 2016. Se busca evitar la fragmentación jurídica entre los estados miembros. Se publica este tratado que no requiere transposición. Aunque los estados pueden elaborar leyes aclarando o desarrollando aspectos. En españa debe ser una Ley Orgánica, por ser la protección de las personas físicas un derecho fundamental. En caso de conflicto entre el Reglamento y la Ley Orgánica, prima el reglamento.

        Su ámbito de aplicación es al tratamiento total o parcialente automatizado de datos personales, así como al tratamiento no automatizado de de datos personales contenidos o destinados a ser incluidos en un fichero. Excepto: actividad no comprendida en el ámbito de aplicación de derecho de la unión, efectuado por personas físicas en el ámbito doméstico o por parte de las autoridades competentes con fines de prevención, investigación, detección o enjuiciamiento penal.

        Se aplica regionalmente en el contexto de un establecimiento del responsable o encargado. Independientemente de que el tratamiento tenga lugar o no en la Unión

    \subsubsection{Definiciones}
        \begin{description}
            \item[Limitación del tratamiento] El marcado de los datos de caracter personal con el fin de limitar su tratamiento en el futuro.
            \item[Datos personales] Toda información sobre una persona física, identificada o identificable.
            \item[Persona identificable] Persona cuya identidad pueda determinarse directa o indirectamente.
            \item[Elaboración de perfiles] toda forma de tratamiento automatizado consistente en usar los datos personales para evaluar determinados aspectos persona;les de una persona física.
            \item[Seudonimización] Proceso por el que los datos ya no pueden atribuirs a una persona sin información adicional.
            \item[Anonimización] Tratamiento por el que se disocian los datos haciéndo imposible ser atribuidos a una persona
            \item[Responsable de tratamiento] La persona, autoridad, servicio u organismo qu determine los medios y fines del tratamiento.
            \item[Encargado de tratamiento] La persona, autoridad, servicio u organismo que trate los datos personales por cuenta del responsable del tratamiento.     
        \end{description}

    \subsubsection{Principios}
        El RGPD señala un conjunto de principios que los responsables y encargados deben observar al tratar los datos personales.
        \begin{description}
            \item[Licitud Transparencia y Lealtad]
            \item[Limitación de la finalidad] Implica la obligatoriedad de que los datos sean tratados con una serie de finalidades explícitas y legítimas, y se prohibe que sean tratados de una manera incompatible con esos fines.
            \item[Minimización de datos] Que los datos sean adecuados, pertinentes y limitados a lo necesario. 
            \item[Exactitud] Deben adoptarse las medidas razonables para que se rectifiquen o supriman datos inexactos.
            \item[Limitación del plazo de conservación] Debe limitarse en el tiempo al logro de los fines que el tratamiento persigue. Una vez terminado el fin, los datos deben ser borrados o anonimizados.
            \item[Integridad y confidencialidad] Debiendo actuar proactivamente en la protección de riesgos que amenacen su seguridad.
            \item[Responsabilidad proactiva] Los responsables deben aplicar las medidas técnicas y organizativas apropiadas para garantizar y demostrar que el tratamiento se lleva a cabo en conformidad con el reglamento.
        \end{description}

    \subsubsection{Licitud del tratamiento}
        El tratamiento es lícito si:
        \begin{itemize}
            \item Hay consentimiento
            \item Es necesario para cumplir un contrato
            \item Hay obligación legisl
            \item Necesario para intereses vitales
            \item Hay interés público
            \item Es necesario para la satisfacción de intereses legítimos perseguidos por el responsable del tratamiento o por un tercero.
        \end{itemize}

    \subsubsection{Categorías de datos especiales}
        Existen datos de especial protección. Son el orígen étnico o racial, las opiniones políticas, las convicciones religiosas o filosóficas, la afiliación sindical, el tratamiento de datos genéticos, los datos biométricos dirigidos a identificar de manera unívoca a una persona, los datos relativos a la salud, la vida sexual o la orientación. Para todos estos datos queda prohibido el tratamiento, excepto si se da una de las siguientes situaciones:
        \begin{itemize}
            \item Hay consentimiento explicito
            \item Es necesario para el cumplimiento de obligaciones y ejercicio de derechos
            \item Para proteger intereses vitales
            \item Asociaciones, partidos políticos, sindicatos
            \item Si el interesado lo ha hecho manifiestamente público
            \item Interés público esencial
        \end{itemize}

    \subsubsection{Derechos en el RGPD}
        El RGPD otorga una serie de derechos a los interesados>
        \begin{itemize}
            \item Transparencia. Los interesados pueden solicitar información y el responsable tendrá que responder en un mes, esta información debe ser gratuita a no ser que sea excesivamente repetitiva o haya generado costes administrativos.
            \item Acceso
            \item Rectificación
            \item Olvido
            \item Limitación del tratamiento
            \item Portabilidad
            \item Oposición de características
            \item no ser objeto de decisiones basadas en el tratamiento automatizado.
        \end{itemize}

    \subsubsection{Obligaciones}
        \paragraph{Obligaciones de responsabilidad}
            El responsable del tratamiento tiene la obligación de cumplir el reglamento y de demostrar su observancia. Se quiere que la protección de datos ocurra desde el diseño y por defecto.

            Cada responsable debe llevar un registro de las actividades del tratamiento. Quedan exentas las empresas de menos de 250 personas, con datos cuyo tratamiento no implique riesgos para derechos y libertades de los interesados, no sea ocasional o incluya datos especiales.

            El responsable y encargado deben cooperar con la autoridad que lo solicite en el desempeño de sus funciones.

        \paragraph{Seguridad de los datos personales}
            El responsable deberá aplicar medidas técnicas y organizarivas para garantizar el nivel de seguridad adecuado al riesgo. En caso de violación, el responsable deberá notificarlo a la autoridad de control y, de ser posible, antes de 72 horas de tener constancia de ello, si supone un riesgo para las libertades.

        \paragraph{Evaluación de impacto}
            Es una metodología de evaluación de riesgos y adopción de medidas necesarias para la protección de datos, evitando y minimizando el impacto. Es necesariollevarlo a cabo en los siguientes tratamientos.

            El responsable deberá realizar una evaluación de impacto cuando sea posible que un tipo de tratamiento, (especialmente con nuevas tecnologías), por su naturaleza, alcance, contexto o fines, entrañe un alto riego. 

            Debe contener descripción de las operaciones de tratamiento previstas, junto con los fines. También las medidas para afrontar riesgos y mecanismos que garanticen la protección de datos personales.

        \paragraph{Análisis de riesgo}
            Para determinar las medidas a aplicar es necesario una evaluación de riesgo. Una vez hecho se podrán tomar las medidas pertinentes.

            La ley obliga a los responsables a aplicar a los tratamientos de datos, las medidas de seguridad que correspondan en el Esquema Nacional de Seguridad, e impulsar un grado de implementación de medidas equivalentes en el ámbito privado.
        
        \paragraph{Transferencias internacionales}
            Se puede realizar una transferencia de datos a un tercer país u organización internacional, cuando la Comisión haya decidido que garantizan un nivel de protección adecuado. Dicha transferencias no requieren de autorización previa. Cuando el destino no esté en las listas de la UE, será posible si el responsable ofrece las garantías adecuadas.

    \subsubsection{Impacto del RGPD en las administraciones públicas}
        El \textbf{Esquema Nacional de Seguridad} y el \textbf{RGPD} establecen la obligación de las administraciones públicas de realizar análisis de riesgo para determinar el impacto de los tratamientos de datos en los derechos y libertades de las personas.

        La \textbf{AEPD} ha publicado un documento en el que pone de manifiesto estas medidas de seguridad.

    \subsubsection{ASSI-RGPD}
        Es una aplicación que permite a los responsables del tratamiento de datos la ejecución de actividades
        \begin{description}
            \item[Riegos] Análisis y evaluación de impacto
            \item[Información] Para que cada TDP que hay que incluir en el registro de actividades de Tratamiento
            \item[Verificación] del resto de aspectos normativos del RGPD
        \end{description}
        A través de una serie de formularios, ofrece un conjunto de informes y documentos que ayudan al cumplimiento de las obligaciones.

    \subsubsection{Sanciones}
        \begin{description}
            \item[Falta Grave] Multa administrativa de 10.000.000€ o 2\% del volumen de negocio anual, lo que sea mayor.
            \item[Falta muy grave] Multa administrativa de 20.000.000€ o 4\% del volumen de negocio anual, lo que sea mayor.
        \end{description}

    \subsubsection{Infoprmación adicional}
    \hyperref{Novedades en materia de protección de datos}{http://www.inap.es/mediateca?p_p_id=contentviewerservice_WAR_alfresco_packportlet&p_p_lifecycle=0&p_p_state=maximized&p_p_mode=view&p_p_col_id=columna-1&p_p_col_pos=1&p_p_col_count=2&_contentviewerservice_WAR_alfresco_packportlet_nodeName=JORNADA_NOVEDADES_PROTECCION_DATOS_2753529.gcl&_contentviewerservice_WAR_alfresco_packportlet_struts_action=\%2Fcontentviewer\%2Fview&contentType=notice} en INAP.

    También hay iunformación adicional en la AEPD.

\subsection{Ley orgánica 3/2018 de protección de datos y garantía de derechos digitales}
        Implementa y expande la RGPD
    \subsubsection{Aspectos mas importantes}
        Se regulan los datos de las \textbf{personas fallecidas} permitiendo que familiares y herededos puedan solicitar el acceso, rectificación y supresión, si así lo instruyó el fallecido.

        Facilita la ejecución de derechos al exigir que los medios sean accesibles. También se regula el modo en que debe informarse a las personas acerca del tratamiento de datos. Facilitando los aspectos importantes y relegando a un enlace la siguiente información.

        Se fija en \textbf{14 años el consentimiento autónomo}, regulando expresamente el derecho de supresión en redes sociales u otros servicios durante la minoría de edad del menor.

        Se regula también el derecho al olvido en redes sociales.

        Se recogen sistemas de denuncias anónimas, para poner en concocimento de entidades privadas, la comisión de actos contrarios a la norma.

        Sen articulan garantías del derecho a la intimidad frente a los dispositivos de videovigilancia en el trabajo. Complementando con los sistemas de geolocalización en el ámbito laboral, de los que deberán estar informados.

        Se regulan los sistemas de información crediticia, reduciendo a 5 años el periodo de inclusión de las deudas y teniendo un precio mínimo de 50€.

        Se modifica la ley de competencia desleal, regulando como prácticas agresivas la suplantación de la Agencia o sus funciones.

        En la constitución se establece el mandato de la garantía de los derechos digitales, proponiéndose esto en la LOPD

    \subsubsection{Garantía de derechos digitales}
        La ley recoge una serie de derechos y garantías que se pueden organizar en bloques:
        \begin{description}
            \item[Internet] Derecho a la neutralidad de la red, protección de menores, acceso universal, olvido y rectificación.
            \item[Digital] Seguridad digiral, actualización de la información, educación digital y testamento.
            \item[Laboral] Derechos digitales en la negociación colectiva, desconexión digiral e intimidad en el uso de dispositivo en la geolocalización
            \item[Redes sociales] Olvido y portabilidad
        \end{description}
    \subsubsection{Prescripción de las sanciones}
        Se establece una serie de plazos según el importe de la cuantía:
        \begin{table}[H]
        \begin{tabular}{rl}
            hasta 40.000€ & prescriben al año \\
            hasta 300.000€ & prescriben a los dos años \\
            mas de 300.000€ & Prescriben a los 3 años \\
        \end{tabular}
        \end{table}

    \subsubsection{Agencia Española de Protección de Datos}
        Regula las funciones y potestados de la AEPD como \textbf{autoridad administrativa independiente}: Supervisa la aplicación de esta ley orgánica y del Reglamento General. Queda fuera de la ley los derechos en la era digital que se asignarán a otro organismo no determinado. 

    \subsubsection{Herramientas de la AEPD}
        Ha desarrollado una herramienta llamada Facilita RGPD para profesionales. Permite valorar la situación con respecto al tratamiento de datos personales. Determinando si quedan acciones pendientes para cumplir el reglamento.

    \subsubsection{Protección de datos y procedimiento administrativo común}
        Se modifica la ley del Procedimiento Administrativo Común, indicando que no aportar un documento es un derecho del interesado (en lugar de no estar obligados). Independientemente de si es o no la administración actuante, deberán recabarlos usando sus redes corporativas, plataforma de mediación u otros sistemas a tal efecto, salvo que el interesado se oponga a ello (que no podrá en caso de procedimientos sancionadores o de inspección).

        En la \textbf{disposición adicional séptima} se indica sobre las notificaciones por medio de anuncios y publicaciones de actos administrativos.

        Se anonimizarán los datos de los ciudadanos con nombre y apellidos y cuatro cifras aleatorias de su NIF (si lo hubiere). Nunca debe publicarse conjuntamente nombre apellidos y NIF

        Se elaborará un protocolo específico de publicación de datos de violencia de género.

        Las administraciones públicas tienen la potestad deverificación. Cuando se formulan solicitudes en las que el interesado declare datos personales, el órgano destinatario podrá efectuar las verificaciones necesarias para comprobar la exactitud de los datos.


    \subsubsection{El delegado de protección de datos en la administración pública}
        Es la figura responsable de la privacidad, con función proactiva y preventiva. Tiene que tener conocimientos especializados de derecho. 
        
        Es obligatorio designar un RPD si:
        \begin{itemize}
            \item Administraciones públicas
            \item Tratamiento de datos a gran especialización
            \item Tratamiento de datos especiales o relativos a condenas e infracciones penales
        \end{itemize}
        
        El delegado tendrá como mínimo las siguientes funciones:
        \begin{itemize}
            \item Informar y asesorar al responsable o encargado de tratamiento, y a los empleados que se ocumen de estos.
            \item Supervisar el cumplimiento en el Reglamento, o de otras disposiciones
            \item Supervisar la asignación de responsabilidades
            \item Supervisar la concienciación y formación del personal que participa en las operaciones
            \item Supervisar auditorías
            \item Ofrecer asesoramiento sobre la evaluación de impacto
            \item Supervisar su aplicación de conformidad con el artículo 35
            \item Cooperar con autoridad de control
            \item Actuar como punto de contacto de la autoridad de control para cuestiones relativas al tratamiento.
            \item Realizar consiltas a la autoridad de control, sobre cualquier otro asunto.
        \end{itemize}

        El nombramiento y cese debe comunicarse a la AEPD.

